% GENERATED FILE. Edit canonical lessons, then run npm run generate:books.
% canonical-source-hash: fnv1a64:648cd7804f278be6
% canonical-lessons: FR-C02-je, FR-C02-me, FR-C02-appeler, FR-C02-je-mappelle, FR-C02-tu-vous, FR-C02-comment, FR-C02-comment-vous-appelez-vous, FR-C02-enchante, FR-C02-practice

\chapter{Introducing Yourself}
\label{ch:introductions}
\hlchaptermodality{2}

\begin{chapteropening}
\textbf{By the end of this chapter:} \emph{I can introduce myself in French, give my name, and understand the other person's name.}

Now a first real exchange — “Je m'appelle Mira.” / “Comment vous appelez-vous?” / “Je m'appelle Arun.” / “Enchanté.” — giving your name, asking for someone else's, and saying you are glad to meet them. We build it the same way as the greetings: each piece taken apart first, then assembled.
\end{chapteropening}

\section[je]{je — \textquotedblleft{}I,\textquotedblright{} the smallest word you own}

\label{lesson:FR-C02-je}



\begin{quote}
To introduce yourself you need \textquotedblleft{}I.\textquotedblright{} French's is tiny — but it comes from a word you already own whole in English.
\end{quote}

\begin{sounds}
\begin{itemize}
\raggedright
  \item \textbf{je} = \emph{zhuh} — a soft \emph{j} (as in \emph{measure}) plus a light \emph{uh}.
  \item Before a vowel it \textbf{elides}: \emph{je ai} $\to$ \textbf{j'ai} (\textquotedblleft{}I have\textquotedblright{}). You'll see the same squeeze in \emph{je m'appelle}.
\end{itemize}
\end{sounds}

\begin{cousinweb}
\textbf{je} — \textquotedblleft{}I.\textquotedblright{} Root: Latin \textbf{ego}. That Latin form survives \emph{raw} in English:

\begin{itemize}
\raggedright
  \item \textbf{ego}, \textbf{ego}ist, \textbf{ego}tism, \textbf{ego}centric — the \textquotedblleft{}I\textquotedblright{} at the center of the self.
\end{itemize}

So the word for \textquotedblleft{}I\textquotedblright{} and the psychologist's \emph{ego} are the \textbf{same} ancient word — one worn down to a whisper (\emph{je}), one kept whole (\emph{ego}). (Spanish wore \emph{ego} down to \emph{yo}; French wore it further, to \emph{je}.)
\end{cousinweb}

\subsection*{Your turn}
Say these aloud:

\begin{itemize}
\raggedright
  \item \textquotedblleft{}je\textquotedblright{} (\emph{zhuh})
  \item the elision — \textquotedblleft{}je ai\textquotedblright{} $\to$ \textquotedblleft{}j'ai\textquotedblright{}
  \item the English cousin that kept the Latin whole — \emph{ego}
\end{itemize}

\begin{tcolorbox}[breakable,colback=teal!4,colframe=teal!35!black,title={Before you move on}]
What Latin word is \textbf{je} worn down from, and what English word kept it whole? (\emph{ego} — English \emph{ego}.) What happens to \emph{je} before a vowel? (It elides to \emph{j'}.)
\end{tcolorbox}
\section[me]{me — \textquotedblleft{}myself,\textquotedblright{} the same word as English \emph{me}}

\label{lesson:FR-C02-me}



\begin{quote}
The middle piece of \emph{je m'appelle}. It's a real word with a real history — not just a bit of glue — and you've owned it in English all your life.
\end{quote}

\begin{sounds}
\begin{itemize}
\raggedright
  \item \textbf{me} = \emph{muh} — a light \emph{uh}, like \emph{je}.
  \item Before a vowel it \textbf{elides}: \emph{me appelle} $\to$ \textbf{m'appelle} (the \emph{e} drops, the words fuse).
\end{itemize}
\end{sounds}

\begin{cousinweb}
\textbf{me} — \textquotedblleft{}me / myself.\textquotedblright{} Root: Latin \textbf{mē} (the accusative \textquotedblleft{}me\textquotedblright{}) — literally the \textbf{same} word as English \textbf{me}. The whole family grows from one Proto-Indo-European root \textbf{*me-}:

\begin{itemize}
\raggedright
  \item English \textbf{me}, \textbf{my}, \textbf{mine};
  \item Latin \textbf{mē} and \textbf{meus} (\textquotedblleft{}my\textquotedblright{}), which became French \emph{mon / ma / mes}.
\end{itemize}

So the tiny \emph{m'} in \emph{je m'appelle} is a first-person marker you've carried in English since childhood — \emph{me}, \emph{my}, \emph{mine} are its cousins.
\end{cousinweb}

\begin{grammarlens}[title={Grammar Lens: the reflexive set (me / te / se)}]
\emph{me} is one of three little pronouns that \textbf{bounce a verb back onto its doer} — \emph{me} (\textquotedblleft{}myself\textquotedblright{}), \emph{te} (\textquotedblleft{}yourself\textquotedblright{}), \emph{se} (\textquotedblleft{}himself, herself\textquotedblright{}). Their roots run in parallel with the one you just met, and two of them you already own in English:

\begin{itemize}
\raggedright
  \item \emph{te} $\leftarrow$ Latin \textbf{tē}, the cousin of archaic English \textbf{thee}.
  \item \emph{se} $\leftarrow$ Latin \textbf{sē}, the \textquotedblleft{}self\textquotedblright{} root behind \textbf{se}parate and, at a distance, \textbf{su}icide (\emph{suī}, \textquotedblleft{}of oneself\textquotedblright{}).
\end{itemize}

Only \textbf{me} does any work here: it is what gives your own name, \emph{je m'appelle}. \emph{te} turns up once, in the informal \textquotedblleft{}what's your name?\textquotedblright{} (\emph{comment tu t'appelles?}), and \emph{se} waits for a lesson that needs it. The set is named so its shape is familiar, not laid out as a grid to memorize — you meet each one where it earns its place.
\end{grammarlens}

\subsection*{Your turn}
Say these aloud:

\begin{itemize}
\raggedright
  \item \textquotedblleft{}me\textquotedblright{} (\emph{muh}), then its elision — \textquotedblleft{}m'appelle\textquotedblright{}
  \item the English cousins from the same root — \emph{me}, \emph{my}, \emph{mine}
  \item the set — \emph{me} (myself), \emph{te} (yourself), \emph{se} (oneself)
\end{itemize}

\begin{tcolorbox}[breakable,colback=teal!4,colframe=teal!35!black,title={Before you move on}]
What Latin word is \textbf{me} from, and what three English words share its root? (\emph{mē} — \emph{me}, \emph{my}, \emph{mine}.) What are its two siblings, and what does \emph{me} do to a verb? (\emph{te}, \emph{se}; it bounces the action back onto the doer — reflexive.)
\end{tcolorbox}
\section[(s')appeler]{(s')appeler — \textquotedblleft{}to call,\textquotedblright{} and \textquotedblleft{}to call \emph{oneself}\textquotedblright{}}

\label{lesson:FR-C02-appeler}



\begin{quote}
French says \textquotedblleft{}my name is\textquotedblright{} by saying \textquotedblleft{}I call myself.\textquotedblright{} Here's the verb that does it — and a fat English family riding on it.
\end{quote}

\begin{sounds}
\begin{itemize}
\raggedright
  \item \textbf{appeler} = \emph{a-play} — the \emph{-er} infinitive ending is a plain \textquotedblleft{}ay.\textquotedblright{}
  \item In \textbf{m'appelle} the doubled \emph{ll} keeps a crisp \textquotedblleft{}l\textquotedblright{}: \emph{ma-PELL}.
\end{itemize}
\end{sounds}

\begin{cousinweb}
\textbf{appeler} — \textquotedblleft{}to call.\textquotedblright{} Root: Latin \textbf{appellāre} (\textquotedblleft{}to address, call by name\textquotedblright{}). The English cousins are everywhere in formal language:

\begin{itemize}
\raggedright
  \item \textbf{appeal} — to \emph{call} to a higher authority.
  \item \textbf{appell}ation — a name or title (an \emph{appellation} of wine is its official name).
  \item \textbf{appell}ate, \textbf{appell}ant — the court that hears appeals, and the one who brings it.
\end{itemize}

To \textbf{name} someone is, at root, to \textbf{call} them.
\end{cousinweb}

\begin{grammarlens}[title={Grammar Lens: reflexive verbs — doing something to yourself}]
French builds \textquotedblleft{}to be called\textquotedblright{} from \textquotedblleft{}to call\textquotedblright{} plus a little word that bounces the action back onto the doer: \textbf{s'appeler}, \textquotedblleft{}to call \emph{oneself}.\textquotedblright{} That bounce-back word changes for each person:

\begin{itemize}
\raggedright
  \item \textbf{me} = myself · \textbf{te} = yourself · \textbf{se} = himself/herself.
\end{itemize}

So \textquotedblleft{}I call myself\textquotedblright{} needs \textbf{me}: \emph{je \textbf{me} appelle} $\to$ (elided) \emph{je m'appelle}. English hides the same move in \textquotedblleft{}I call \textbf{myself}\textquotedblright{} — French just uses it as the normal way to give a name.
\end{grammarlens}

\subsection*{Your turn}
Say these aloud:

\begin{itemize}
\raggedright
  \item \textquotedblleft{}appeler\textquotedblright{} (\emph{a-play})
  \item the English cousins — \emph{appeal}, \emph{appellation}
  \item \textquotedblleft{}call oneself\textquotedblright{} with \emph{me} — \emph{je me appelle}
\end{itemize}

\begin{tcolorbox}[breakable,colback=teal!4,colframe=teal!35!black,title={Before you move on}]
What Latin verb is \textbf{appeler} from, and two English cousins? (\emph{appellāre} — \emph{appeal}, \emph{appellation}.) What does a \emph{reflexive} verb do, and which little word means \textquotedblleft{}myself\textquotedblright{}? (Bounces the action back on the doer; \emph{me}.)
\end{tcolorbox}
\section[je m'appelle…]{je m'appelle… — \textquotedblleft{}my name is…,\textquotedblright{} literally \textquotedblleft{}I call myself\textquotedblright{}}

\label{lesson:FR-C02-je-mappelle}



\begin{quote}
Assemble the two pieces you just met — \emph{je} and \emph{(s')appeler} — into the single most useful sentence for meeting anyone.
\end{quote}

\begin{cousinweb}
Three atoms, each already taken apart on its own — je (I, $\leftarrow$ \emph{ego}), me (myself, $\leftarrow$ \emph{mē}), appelle (call, $\leftarrow$ \emph{appellāre}):

\begin{itemize}
\raggedright
  \item \textbf{je} (I) + \textbf{me} (myself) + \textbf{appelle} (call).
  \item \emph{me} elides before the vowel of \emph{appelle} $\to$ \textbf{m'}.
\end{itemize}

\begin{quote}
\textbf{je m'appelle…} = word for word \textquotedblleft{}I call myself…\textquotedblright{} = \textbf{\textquotedblleft{}my name is…\textquotedblright{}}
\end{quote}

\emph{Je m'appelle Arun.} $\to$ \textquotedblleft{}My name is Arun.\textquotedblright{} \emph{Je m'appelle Mira.} $\to$ \textquotedblleft{}My name is Mira.\textquotedblright{}
\end{cousinweb}

\begin{culture}
French \emph{prefers} this verb-based \textquotedblleft{}I call myself\textquotedblright{} to a literal \textquotedblleft{}my name is.\textquotedblright{} The literal form does exist:

\begin{itemize}
\raggedright
  \item \textbf{mon nom est…} (\textquotedblleft{}my name is…\textquotedblright{}), with \textbf{nom} (\textquotedblleft{}name\textquotedblright{}) $\leftarrow$ Latin \textbf{nōmen} $\to$ English \textbf{noun}, \textbf{nom}inate, \textbf{nom}inal, \textbf{nom}enclature.
\end{itemize}

But \emph{mon nom est} sounds stiff and official — a passport-desk phrase. For a real introduction, reach for \textbf{je m'appelle}. (You now also see why \textquotedblleft{}noun\textquotedblright{} is literally \textquotedblleft{}a naming word.\textquotedblright{})
\end{culture}

\subsection*{Your turn}
Say these aloud:

\begin{itemize}
\raggedright
  \item \textquotedblleft{}je m'appelle…\textquotedblright{} then your own name
  \item what it \emph{literally} says (\textquotedblleft{}I call myself\textquotedblright{})
  \item the stiff literal alternative — \textquotedblleft{}mon nom est…\textquotedblright{}
\end{itemize}

\begin{tcolorbox}[breakable,colback=teal!4,colframe=teal!35!black,title={Before you move on}]
What does \textbf{je m'appelle} literally mean, and why does \emph{me} become \emph{m'}? (\textquotedblleft{}I call myself\textquotedblright{}; elision before the vowel of \emph{appelle}.) What's the more literal but stiffer \textquotedblleft{}my name is,\textquotedblright{} and what English word is \emph{nom} cousin to? (\emph{mon nom est}; \emph{noun}.)
\end{tcolorbox}
\section[tu / vous]{tu / vous — \textquotedblleft{}you,\textquotedblright{} familiar and formal}

\label{lesson:FR-C02-tu-vous}



\begin{quote}
To \emph{ask} someone their name, you need \textquotedblleft{}you.\textquotedblright{} French has two, and choosing between them is a small act of social judgment.
\end{quote}

\begin{sounds}
\begin{itemize}
\raggedright
  \item \textbf{tu} = \emph{too}, but with the front-rounded French \emph{u} (lips rounded, tongue forward). \textbf{vous} = \emph{voo} (silent \emph{s}).
\end{itemize}
\end{sounds}

\begin{cousinweb}
\begin{itemize}
\raggedright
  \item \textbf{tu} (familiar \textquotedblleft{}you\textquotedblright{}) $\leftarrow$ Latin \textbf{tū}.
  \item \textbf{vous} (formal / plural \textquotedblleft{}you\textquotedblright{}) $\leftarrow$ Latin \textbf{vōs}, \textquotedblleft{}you all.\textquotedblright{}
\end{itemize}

English once had the same split — old \textbf{thou} was the \emph{tu} — but \emph{thou} died out, leaving one flat \textbf{you} for everyone.
\end{cousinweb}

\begin{grammarlens}[title={Grammar Lens: two \textquotedblleft{}you\textquotedblright{}s, chosen by register}]
\begin{itemize}
\raggedright
  \item \textbf{tu} — for friends, family, children, peers. Intimate, equal.
  \item \textbf{vous} — for strangers, elders, anyone you'd show respect — \textbf{and} it's the plural \textquotedblleft{}you all.\textquotedblright{}
\end{itemize}

The trick: French makes politeness by using the \textbf{plural} on a single person. Addressing one respected person as \textquotedblleft{}you all\textquotedblright{} is the polite move (French even has a verb for it, \emph{vouvoyer}). This is the same politeness Spanish builds with \emph{usted}, but by a different route: Spanish coined a whole new pronoun (from \textquotedblleft{}your grace\textquotedblright{}); French just borrows its plural. Using \emph{tu} too early can feel over-familiar — \textbf{when in doubt, \emph{vous}}.
\end{grammarlens}

\subsection*{Your turn}
Say these aloud:

\begin{itemize}
\raggedright
  \item \textquotedblleft{}tu\textquotedblright{} (front-rounded \emph{u}), \textquotedblleft{}vous\textquotedblright{} (\emph{voo})
  \item who gets \emph{tu} (friends), who gets \emph{vous} (strangers, elders)
  \item the rule of thumb — \textquotedblleft{}when in doubt, \emph{vous}\textquotedblright{}
\end{itemize}

\begin{tcolorbox}[breakable,colback=teal!4,colframe=teal!35!black,title={Before you move on}]
What decides \emph{tu} vs \emph{vous} — and what does French do to be polite to one person? (Register — familiar vs respect; it uses the \emph{plural} (\emph{vous}) on one person.) What old English word was the lost \emph{tu}? (\emph{thou}.)
\end{tcolorbox}
\section[comment]{comment — \textquotedblleft{}how,\textquotedblright{} from the same \emph{quomodo} as Spanish \emph{cómo}}

\label{lesson:FR-C02-comment}



\begin{quote}
French asks \textquotedblleft{}what's your name?\textquotedblright{} as \textquotedblleft{}\textbf{how} are you called?\textquotedblright{} — so first, the word for \textquotedblleft{}how.\textquotedblright{}
\end{quote}

\begin{sounds}
\begin{itemize}
\raggedright
  \item \textbf{comment} = \emph{ko-MAHN} — the final \emph{-ent} is a \textbf{nasal} vowel; the \emph{t} is silent.
\end{itemize}
\end{sounds}

\begin{cousinweb}
\textbf{comment} — \textquotedblleft{}how.\textquotedblright{} Root: Latin \textbf{quo modo}, \textquotedblleft{}in what manner\textquotedblright{} — two words fused and worn into one.

\begin{itemize}
\raggedright
  \item That same \textbf{quomodo} also became Spanish \textbf{cómo} — two Romance languages, one Latin \textquotedblleft{}in-what-way,\textquotedblright{} reshaped differently.
  \item The piece \textbf{modo} (\textquotedblleft{}manner, measure\textquotedblright{}) is alive in English \textbf{mode}, \textbf{mod}el, \textbf{mod}erate, \textbf{mod}ify — all about the \emph{way} a thing is done.
\end{itemize}
\end{cousinweb}

\begin{grammarlens}[title={Grammar Lens: French counts names as a matter of \textquotedblleft{}how\textquotedblright{}}]
Where English asks \textquotedblleft{}\textbf{what} is your name?\textquotedblright{}, French asks \textquotedblleft{}\textbf{how} are you called?\textquotedblright{} — \emph{comment} (how), not \emph{what}. The question is about the \emph{manner} of your calling, not a thing to be named. Don't translate word for word; learn the French logic: the name-question runs on \textbf{comment}.
\end{grammarlens}

\subsection*{Your turn}
Say these aloud:

\begin{itemize}
\raggedright
  \item \textquotedblleft{}comment\textquotedblright{} (\emph{ko-MAHN}, nasal ending)
  \item the Latin \emph{quo modo} $\to$ French \emph{comment}, Spanish \emph{cómo}
  \item English \emph{mode/model} share the \emph{modo} piece
\end{itemize}

\begin{tcolorbox}[breakable,colback=teal!4,colframe=teal!35!black,title={Before you move on}]
What two Latin words fused into \textbf{comment}, and what Spanish word shares them? (\emph{quo modo} — Spanish \emph{cómo}.) English asks \textquotedblleft{}what is your name?\textquotedblright{} — what does French literally ask? (\textquotedblleft{}\emph{How} are you called?\textquotedblright{})
\end{tcolorbox}
\section[comment vous appelez-vous?]{comment vous appelez-vous? — \textquotedblleft{}what's your name?\textquotedblright{}}

\label{lesson:FR-C02-comment-vous-appelez-vous}



\begin{quote}
Assemble it: \textquotedblleft{}how\textquotedblright{} + \textquotedblleft{}do you call yourself\textquotedblright{} = the polite question that gets the introduction going.
\end{quote}

\begin{cousinweb}
\begin{itemize}
\raggedright
  \item \textbf{comment} (how) + \textbf{vous appelez-vous} (\textquotedblleft{}do you call yourself\textquotedblright{}).
\end{itemize}

\begin{quote}
\textbf{comment vous appelez-vous?} = literally \textquotedblleft{}\textbf{how do you call yourself?}\textquotedblright{} = the polite \textbf{\textquotedblleft{}what's your name?\textquotedblright{}}
\end{quote}
\end{cousinweb}

\begin{grammarlens}[title={Grammar Lens: asking by inversion}]
A statement — \emph{vous vous appelez} (\textquotedblleft{}you call yourself\textquotedblright{}) — becomes a \textbf{question} by flipping the verb and the pronoun, joined with a hyphen:

\begin{quote}
\emph{vous appelez-\textbf{vous}?}
\end{quote}

That flip is French's formal way to ask anything. The \textbf{informal} version simply swaps in \emph{tu}:

\begin{quote}
\emph{comment \textbf{tu t'appelles}?} — note \emph{te} elides to \textbf{t'}, exactly as \emph{me} became \emph{m'} in \emph{je m'appelle}.
\end{quote}

So the same verb, \emph{s'appeler}, runs the whole exchange: \emph{je m'appelle} to give your name, \emph{comment vous appelez-vous?} to ask for theirs.
\end{grammarlens}

\subsection*{Your turn}
Say these aloud:

\begin{itemize}
\raggedright
  \item \textquotedblleft{}comment vous appelez-vous?\textquotedblright{} (formal)
  \item the informal — \textquotedblleft{}comment tu t'appelles?\textquotedblright{}
  \item the pair — ask, then answer: \textquotedblleft{}…vous appelez-vous?\textquotedblright{} / \textquotedblleft{}je m'appelle…\textquotedblright{}
\end{itemize}

\begin{tcolorbox}[breakable,colback=teal!4,colframe=teal!35!black,title={Before you move on}]
What does \textbf{comment vous appelez-vous?} literally say? (\textquotedblleft{}How do you call yourself?\textquotedblright{}) How do you turn the statement into a question, and what's the \emph{tu} version? (Invert verb + pronoun with a hyphen; \emph{comment tu t'appelles?})
\end{tcolorbox}
\section[enchanté(e)]{enchanté(e) — \textquotedblleft{}delighted,\textquotedblright{} i.e. \textquotedblleft{}pleased to meet you\textquotedblright{}}

\label{lesson:FR-C02-enchante}



\begin{quote}
The graceful close to an introduction — and it literally says you've been put under a spell.
\end{quote}

\begin{sounds}
\begin{itemize}
\raggedright
  \item \textbf{enchanté} = \emph{ahn-shahn-TAY} — two \textbf{nasal} vowels (\emph{en}, \emph{an}), then a clean \textquotedblleft{}tay.\textquotedblright{} \textbf{ch} = English \emph{sh}.
\end{itemize}
\end{sounds}

\begin{cousinweb}
\textbf{enchanté} — \textquotedblleft{}delighted.\textquotedblright{} Root: Latin \textbf{in-} (\textquotedblleft{}upon\textquotedblright{}) + \textbf{cantāre} (\textquotedblleft{}to sing, chant\textquotedblright{}). To \emph{enchant} was literally to \textbf{sing a spell upon} someone. The magic survives in the English family:

\begin{itemize}
\raggedright
  \item \textbf{enchant}, \textbf{incant}ation, and the plain \textbf{chant}.
\end{itemize}

So saying \emph{enchanté} on meeting someone means, gallantly, \textquotedblleft{}I'm \textbf{enchanted} {[}to meet you{]}.\textquotedblright{}
\end{cousinweb}

\begin{grammarlens}[title={Grammar Lens: it agrees with \emph{you}}]
\emph{enchanté} follows the gender-agreement rule you met in \emph{bon/bonne}: it reports the \textbf{speaker's own gender}.

\begin{itemize}
\raggedright
  \item A man says \textbf{enchanté}.
  \item A woman says \textbf{enchantée} — an extra \emph{-e} (same sound, silent).
\end{itemize}
\end{grammarlens}

\subsection*{Your turn}
Say these aloud:

\begin{itemize}
\raggedright
  \item \textquotedblleft{}enchanté\textquotedblright{} (\emph{ahn-shahn-TAY})
  \item the English cousins — \emph{enchant}, \emph{incantation}, \emph{chant}
  \item both forms — \emph{enchanté} (m.) / \emph{enchantée} (f.)
\end{itemize}

\begin{tcolorbox}[breakable,colback=teal!4,colframe=teal!35!black,title={Before you move on}]
What does \textbf{enchanté} literally mean, and from what Latin pieces? (\textquotedblleft{}Enchanted\textquotedblright{} — \emph{in-} + \emph{cantāre}, \textquotedblleft{}to sing upon.\textquotedblright{}) How does a woman say it, and why? (\emph{enchantée} — it agrees with the speaker's gender.)
\end{tcolorbox}
\section[Practice]{Chapter 2 — the introduction, whole}

\label{lesson:FR-C02-practice}



\begin{quote}
Every piece is built. Here's the exchange it was all for.
\end{quote}

\subsection*{The exchange}
\noindent
\begin{tabularx}{\linewidth}{@{}>{\raggedright\arraybackslash}X>{\raggedright\arraybackslash}X@{}}
\toprule
\textbf{French} & \textbf{English} \\
\midrule
\emph{Je m'appelle Mira.} & My name is Mira. \\
\emph{Comment vous appelez-vous?} & What's your name? \\
\emph{Je m'appelle Arun.} & My name is Arun. \\
\emph{Enchanté.} & Pleased to meet you. \\
\bottomrule
\end{tabularx}

Every piece traced to a root:

\begin{itemize}
\raggedright
  \item \textbf{je} $\leftarrow$ \emph{ego} (English \emph{ego})
  \item \textbf{me} $\leftarrow$ \emph{mē} (English \emph{me}, \emph{my}, \emph{mine})
  \item \textbf{(s')appeler} $\leftarrow$ \emph{appellāre} (English \emph{appeal}, \emph{appellation})
  \item \textbf{comment} $\leftarrow$ \emph{quo modo} (Spanish \emph{cómo}, English \emph{mode})
  \item \textbf{enchanté} $\leftarrow$ \emph{in-cantāre} (English \emph{enchant}, \emph{incantation}, \emph{chant})
\end{itemize}

And two habits of French to keep: it gives a name with a \textbf{reflexive verb} (\textquotedblleft{}I call \emph{myself}\textquotedblright{}) and asks for one with \textbf{\textquotedblleft{}how\textquotedblright{}} (\textquotedblleft{}\emph{how} are you called?\textquotedblright{}), and it makes politeness by using the \textbf{plural} \emph{vous} on one person.

\subsection*{Your turn}
Say these aloud:

\begin{itemize}
\raggedright
  \item the whole exchange, both voices
  \item your own introduction — \textquotedblleft{}je m'appelle …, enchanté(e)\textquotedblright{}
  \item ask it back — \textquotedblleft{}comment vous appelez-vous?\textquotedblright{}
\end{itemize}

\begin{tcolorbox}[breakable,colback=teal!4,colframe=teal!35!black,title={Before you move on}]
Give your name in French, ask someone else's, and say you're pleased to meet them. (\emph{Je m'appelle… / Comment vous appelez-vous? / Enchanté(e).}) Which two French habits does this exchange show? (Naming with a reflexive verb; asking with \textquotedblleft{}how,\textquotedblright{} not \textquotedblleft{}what.\textquotedblright{})

Next chapter: \emph{Comment ça va?} (\textquotedblleft{}how's it going?\textquotedblright{}) and answering with \emph{bien} again — the responding cycle.
\end{tcolorbox}
